% Figure: Four Levels of Coordination Complexity
% For inclusion in Section 1.2 (What Is a Multi-Agent System?)
% Levels 1-2 adapted from Cognition (2025), extended with Levels 3-4
%
% Usage: % Figure: Four Levels of Coordination Complexity
% For inclusion in Section 1.2 (What Is a Multi-Agent System?)
% Levels 1-2 adapted from Cognition (2025), extended with Levels 3-4
%
% Usage: % Figure: Four Levels of Coordination Complexity
% For inclusion in Section 1.2 (What Is a Multi-Agent System?)
% Levels 1-2 adapted from Cognition (2025), extended with Levels 3-4
%
% Usage: % Figure: Four Levels of Coordination Complexity
% For inclusion in Section 1.2 (What Is a Multi-Agent System?)
% Levels 1-2 adapted from Cognition (2025), extended with Levels 3-4
%
% Usage: \input{figures/coordination-levels.tex}

\begin{figure*}[t]
\centering
\begin{tikzpicture}[
  agent/.style={draw, thick, rounded corners=3pt, minimum width=1.1cm,
    minimum height=0.8cm, font=\footnotesize\bfseries, inner sep=2pt},
  human/.style={draw, thick, rounded corners=3pt, minimum width=1.1cm,
    minimum height=0.8cm, font=\footnotesize\bfseries, inner sep=2pt,
    fill=orange!12, draw=orange!60!black},
  shared/.style={draw, thick, rounded corners=3pt,
    minimum width=2.4cm, minimum height=0.6cm,
    font=\footnotesize, inner sep=2pt, fill=yellow!10,
    draw=yellow!60!black},
  flow/.style={-{Stealth[length=4pt, width=3pt]}, thick},
  biflow/.style={{Stealth[length=3pt]}-{Stealth[length=3pt]}, thick,
    blue!60!black},
  feedback/.style={{Stealth[length=3pt]}-{Stealth[length=3pt]}, thick,
    red!60!black, dashed},
  panellabel/.style={font=\small\bfseries, anchor=north west},
  levellabel/.style={font=\footnotesize\itshape, text=black!60},
]

% ── Panel positions ──────────────────────────────────────
\def\panelW{7.8}  % panel width
\def\panelH{3.8}  % panel height
\def\gapX{0.6}    % horizontal gap
\def\gapY{0.8}    % vertical gap

% Panel origins (bottom-left of each)
\coordinate (P0) at (0, \panelH + \gapY);
\coordinate (P1) at (\panelW + \gapX, \panelH + \gapY);
\coordinate (P2) at (0, 0);
\coordinate (P3) at (\panelW + \gapX, 0);

% ── Panel 0: Single Agent ───────────────────────────────
\begin{scope}[shift={(P0)}]
  \draw[black!20, rounded corners=6pt] (-0.3, -0.3)
    rectangle (\panelW + 0.3, \panelH + 0.3);
  \node[panellabel] at (0, \panelH + 0.15) {Level 0: Single Agent};
  \node[levellabel] at (\panelW, \panelH + 0.15) [anchor=north east]
    {No coordination};

  \node[agent, fill=blue!12, minimum width=2.2cm, minimum height=1.4cm]
    (a0) at (0.5*\panelW, 0.45*\panelH) {Agent};

  \node[font=\scriptsize, text=black!50, align=center]
    at (0.5*\panelW, 0.45*\panelH - 1.3)
    {All reasoning in one context.\\No dependencies to manage.};
\end{scope}

% ── Panel 1: Sequential Pipeline ────────────────────────
\begin{scope}[shift={(P1)}]
  \draw[black!20, rounded corners=6pt] (-0.3, -0.3)
    rectangle (\panelW + 0.3, \panelH + 0.3);
  \node[panellabel] at (0, \panelH + 0.15) {Level 1: Sequential Pipeline};
  \node[levellabel] at (\panelW, \panelH + 0.15) [anchor=north east]
    {Design-time coordination};

  \node[agent, fill=blue!12] (a1) at (1.2, 0.5*\panelH) {A};
  \node[agent, fill=blue!12] (b1) at (3.2, 0.5*\panelH) {B};
  \node[agent, fill=blue!12] (c1) at (5.2, 0.5*\panelH) {C};
  \node[agent, fill=green!12] (d1) at (7.0, 0.5*\panelH) {Out};

  \draw[flow] (a1) -- (b1);
  \draw[flow] (b1) -- (c1);
  \draw[flow] (c1) -- (d1);

  \node[font=\scriptsize, text=black!50, align=center]
    at (0.5*\panelW, 0.5*\panelH - 1.3)
    {Unidirectional flow. No feedback.\\Human architect resolves all coordination.};
\end{scope}

% ── Panel 2: Parallel Fork-Join ─────────────────────────
\begin{scope}[shift={(P2)}]
  \draw[black!20, rounded corners=6pt] (-0.3, -0.3)
    rectangle (\panelW + 0.3, \panelH + 0.3);
  \node[panellabel] at (0, \panelH + 0.15) {Level 2: Parallel Fork-Join};
  \node[levellabel] at (\panelW, \panelH + 0.15) [anchor=north east]
    {Orchestrator at merge};

  \node[agent, fill=purple!12] (orch2) at (1.0, 0.5*\panelH) {Orch};
  \node[agent, fill=blue!12] (w2a) at (3.5, 0.5*\panelH + 1.0) {B};
  \node[agent, fill=blue!12] (w2b) at (3.5, 0.5*\panelH) {C};
  \node[agent, fill=blue!12] (w2c) at (3.5, 0.5*\panelH - 1.0) {D};
  \node[agent, fill=purple!12] (merge2) at (6.0, 0.5*\panelH) {Merge};
  \node[agent, fill=green!12] (out2) at (7.4, 0.5*\panelH) {Out};

  \draw[flow] (orch2) -- (w2a);
  \draw[flow] (orch2) -- (w2b);
  \draw[flow] (orch2) -- (w2c);
  \draw[flow] (w2a) -- (merge2);
  \draw[flow] (w2b) -- (merge2);
  \draw[flow] (w2c) -- (merge2);
  \draw[flow] (merge2) -- (out2);

  % No inter-worker communication
  \draw[thick, red!40, dashed] (3.5, 0.5*\panelH + 0.5)
    -- node[font=\tiny, text=red!50, right, xshift=2pt] {no comm.}
    (3.5, 0.5*\panelH - 0.5);

  \node[font=\scriptsize, text=black!50, align=center]
    at (0.5*\panelW, -0.05)
    {Concurrent but isolated. Workers don't know each other exist.};
\end{scope}

% ── Panel 3: True MAS (Levels 3-4) ─────────────────────
\begin{scope}[shift={(P3)}]
  \draw[black!30, rounded corners=6pt, line width=1.2pt] (-0.3, -0.3)
    rectangle (\panelW + 0.3, \panelH + 0.3);
  \node[panellabel] at (0, \panelH + 0.15)
    {Levels 3--4: True Multi-Agent System};
  \node[levellabel, text=blue!60!black] at (\panelW, \panelH + 0.15)
    [anchor=north east] {Runtime coordination};

  % Shared state
  \node[shared] (ss) at (0.5*\panelW, 0.5*\panelH + 1.5)
    {Shared State};

  % Agents
  \node[agent, fill=blue!12] (ma) at (1.8, 0.5*\panelH) {A};
  \node[agent, fill=red!10] (mb) at (4.0, 0.5*\panelH + 0.6) {B};
  \node[agent, fill=blue!12] (mc) at (4.0, 0.5*\panelH - 0.6) {C};
  \node[agent, fill=red!10] (md) at (6.2, 0.5*\panelH) {D};

  % Human
  \node[human] (mh) at (0.5*\panelW, 0.5*\panelH - 1.5) {Human};

  % Bidirectional agent communication
  \draw[biflow] (ma) -- (mb);
  \draw[biflow] (ma) -- (mc);
  \draw[biflow] (mb) -- (md);
  \draw[biflow] (mc) -- (md);
  \draw[feedback] (mb) -- (mc);

  % Shared state connections
  \draw[flow, yellow!60!black, thin] (ma.north) -- (ss.south west);
  \draw[flow, yellow!60!black, thin] (mb.north) -- (ss.south);
  \draw[flow, yellow!60!black, thin] (md.north) -- (ss.south east);

  % Human connections
  \draw[biflow, orange!60!black] (mh) -- (ma);
  \draw[biflow, orange!60!black] (mh) -- (md);

  \node[font=\scriptsize, text=black!50, align=center]
    at (0.5*\panelW, -0.05)
    {Feedback loops, negotiation, shared state, self-interest, humans as peers.};
\end{scope}

\end{tikzpicture}
\caption{\textbf{The coordination complexity spectrum.}
  \emph{Level~0}: a single agent handles all reasoning in one context.
  \emph{Level~1}: agents execute sequentially with unidirectional flow; all
  coordination is resolved by the human architect at design time (adapted
  from~\citealt{cognition2025dontbuild}).
  \emph{Level~2}: subagents execute concurrently but cannot communicate with
  each other; the orchestrator handles coordination at merge (adapted
  from~\citealt{cognition2025dontbuild}).
  \emph{Levels~3--4}: true multi-agent systems with runtime dependencies,
  bidirectional communication (blue), feedback loops (red dashed), shared state
  (yellow), self-interested agents (red-tinted), and humans participating as
  peers (orange).  Classical MAS research studied Levels~3--4.  Most current
  LLM agent frameworks operate at Levels~1--2.}
\label{fig:coord-levels}
\end{figure*}


\begin{figure*}[t]
\centering
\begin{tikzpicture}[
  agent/.style={draw, thick, rounded corners=3pt, minimum width=1.1cm,
    minimum height=0.8cm, font=\footnotesize\bfseries, inner sep=2pt},
  human/.style={draw, thick, rounded corners=3pt, minimum width=1.1cm,
    minimum height=0.8cm, font=\footnotesize\bfseries, inner sep=2pt,
    fill=orange!12, draw=orange!60!black},
  shared/.style={draw, thick, rounded corners=3pt,
    minimum width=2.4cm, minimum height=0.6cm,
    font=\footnotesize, inner sep=2pt, fill=yellow!10,
    draw=yellow!60!black},
  flow/.style={-{Stealth[length=4pt, width=3pt]}, thick},
  biflow/.style={{Stealth[length=3pt]}-{Stealth[length=3pt]}, thick,
    blue!60!black},
  feedback/.style={{Stealth[length=3pt]}-{Stealth[length=3pt]}, thick,
    red!60!black, dashed},
  panellabel/.style={font=\small\bfseries, anchor=north west},
  levellabel/.style={font=\footnotesize\itshape, text=black!60},
]

% ── Panel positions ──────────────────────────────────────
\def\panelW{7.8}  % panel width
\def\panelH{3.8}  % panel height
\def\gapX{0.6}    % horizontal gap
\def\gapY{0.8}    % vertical gap

% Panel origins (bottom-left of each)
\coordinate (P0) at (0, \panelH + \gapY);
\coordinate (P1) at (\panelW + \gapX, \panelH + \gapY);
\coordinate (P2) at (0, 0);
\coordinate (P3) at (\panelW + \gapX, 0);

% ── Panel 0: Single Agent ───────────────────────────────
\begin{scope}[shift={(P0)}]
  \draw[black!20, rounded corners=6pt] (-0.3, -0.3)
    rectangle (\panelW + 0.3, \panelH + 0.3);
  \node[panellabel] at (0, \panelH + 0.15) {Level 0: Single Agent};
  \node[levellabel] at (\panelW, \panelH + 0.15) [anchor=north east]
    {No coordination};

  \node[agent, fill=blue!12, minimum width=2.2cm, minimum height=1.4cm]
    (a0) at (0.5*\panelW, 0.45*\panelH) {Agent};

  \node[font=\scriptsize, text=black!50, align=center]
    at (0.5*\panelW, 0.45*\panelH - 1.3)
    {All reasoning in one context.\\No dependencies to manage.};
\end{scope}

% ── Panel 1: Sequential Pipeline ────────────────────────
\begin{scope}[shift={(P1)}]
  \draw[black!20, rounded corners=6pt] (-0.3, -0.3)
    rectangle (\panelW + 0.3, \panelH + 0.3);
  \node[panellabel] at (0, \panelH + 0.15) {Level 1: Sequential Pipeline};
  \node[levellabel] at (\panelW, \panelH + 0.15) [anchor=north east]
    {Design-time coordination};

  \node[agent, fill=blue!12] (a1) at (1.2, 0.5*\panelH) {A};
  \node[agent, fill=blue!12] (b1) at (3.2, 0.5*\panelH) {B};
  \node[agent, fill=blue!12] (c1) at (5.2, 0.5*\panelH) {C};
  \node[agent, fill=green!12] (d1) at (7.0, 0.5*\panelH) {Out};

  \draw[flow] (a1) -- (b1);
  \draw[flow] (b1) -- (c1);
  \draw[flow] (c1) -- (d1);

  \node[font=\scriptsize, text=black!50, align=center]
    at (0.5*\panelW, 0.5*\panelH - 1.3)
    {Unidirectional flow. No feedback.\\Human architect resolves all coordination.};
\end{scope}

% ── Panel 2: Parallel Fork-Join ─────────────────────────
\begin{scope}[shift={(P2)}]
  \draw[black!20, rounded corners=6pt] (-0.3, -0.3)
    rectangle (\panelW + 0.3, \panelH + 0.3);
  \node[panellabel] at (0, \panelH + 0.15) {Level 2: Parallel Fork-Join};
  \node[levellabel] at (\panelW, \panelH + 0.15) [anchor=north east]
    {Orchestrator at merge};

  \node[agent, fill=purple!12] (orch2) at (1.0, 0.5*\panelH) {Orch};
  \node[agent, fill=blue!12] (w2a) at (3.5, 0.5*\panelH + 1.0) {B};
  \node[agent, fill=blue!12] (w2b) at (3.5, 0.5*\panelH) {C};
  \node[agent, fill=blue!12] (w2c) at (3.5, 0.5*\panelH - 1.0) {D};
  \node[agent, fill=purple!12] (merge2) at (6.0, 0.5*\panelH) {Merge};
  \node[agent, fill=green!12] (out2) at (7.4, 0.5*\panelH) {Out};

  \draw[flow] (orch2) -- (w2a);
  \draw[flow] (orch2) -- (w2b);
  \draw[flow] (orch2) -- (w2c);
  \draw[flow] (w2a) -- (merge2);
  \draw[flow] (w2b) -- (merge2);
  \draw[flow] (w2c) -- (merge2);
  \draw[flow] (merge2) -- (out2);

  % No inter-worker communication
  \draw[thick, red!40, dashed] (3.5, 0.5*\panelH + 0.5)
    -- node[font=\tiny, text=red!50, right, xshift=2pt] {no comm.}
    (3.5, 0.5*\panelH - 0.5);

  \node[font=\scriptsize, text=black!50, align=center]
    at (0.5*\panelW, -0.05)
    {Concurrent but isolated. Workers don't know each other exist.};
\end{scope}

% ── Panel 3: True MAS (Levels 3-4) ─────────────────────
\begin{scope}[shift={(P3)}]
  \draw[black!30, rounded corners=6pt, line width=1.2pt] (-0.3, -0.3)
    rectangle (\panelW + 0.3, \panelH + 0.3);
  \node[panellabel] at (0, \panelH + 0.15)
    {Levels 3--4: True Multi-Agent System};
  \node[levellabel, text=blue!60!black] at (\panelW, \panelH + 0.15)
    [anchor=north east] {Runtime coordination};

  % Shared state
  \node[shared] (ss) at (0.5*\panelW, 0.5*\panelH + 1.5)
    {Shared State};

  % Agents
  \node[agent, fill=blue!12] (ma) at (1.8, 0.5*\panelH) {A};
  \node[agent, fill=red!10] (mb) at (4.0, 0.5*\panelH + 0.6) {B};
  \node[agent, fill=blue!12] (mc) at (4.0, 0.5*\panelH - 0.6) {C};
  \node[agent, fill=red!10] (md) at (6.2, 0.5*\panelH) {D};

  % Human
  \node[human] (mh) at (0.5*\panelW, 0.5*\panelH - 1.5) {Human};

  % Bidirectional agent communication
  \draw[biflow] (ma) -- (mb);
  \draw[biflow] (ma) -- (mc);
  \draw[biflow] (mb) -- (md);
  \draw[biflow] (mc) -- (md);
  \draw[feedback] (mb) -- (mc);

  % Shared state connections
  \draw[flow, yellow!60!black, thin] (ma.north) -- (ss.south west);
  \draw[flow, yellow!60!black, thin] (mb.north) -- (ss.south);
  \draw[flow, yellow!60!black, thin] (md.north) -- (ss.south east);

  % Human connections
  \draw[biflow, orange!60!black] (mh) -- (ma);
  \draw[biflow, orange!60!black] (mh) -- (md);

  \node[font=\scriptsize, text=black!50, align=center]
    at (0.5*\panelW, -0.05)
    {Feedback loops, negotiation, shared state, self-interest, humans as peers.};
\end{scope}

\end{tikzpicture}
\caption{\textbf{The coordination complexity spectrum.}
  \emph{Level~0}: a single agent handles all reasoning in one context.
  \emph{Level~1}: agents execute sequentially with unidirectional flow; all
  coordination is resolved by the human architect at design time (adapted
  from~\citealt{cognition2025dontbuild}).
  \emph{Level~2}: subagents execute concurrently but cannot communicate with
  each other; the orchestrator handles coordination at merge (adapted
  from~\citealt{cognition2025dontbuild}).
  \emph{Levels~3--4}: true multi-agent systems with runtime dependencies,
  bidirectional communication (blue), feedback loops (red dashed), shared state
  (yellow), self-interested agents (red-tinted), and humans participating as
  peers (orange).  Classical MAS research studied Levels~3--4.  Most current
  LLM agent frameworks operate at Levels~1--2.}
\label{fig:coord-levels}
\end{figure*}


\begin{figure*}[t]
\centering
\begin{tikzpicture}[
  agent/.style={draw, thick, rounded corners=3pt, minimum width=1.1cm,
    minimum height=0.8cm, font=\footnotesize\bfseries, inner sep=2pt},
  human/.style={draw, thick, rounded corners=3pt, minimum width=1.1cm,
    minimum height=0.8cm, font=\footnotesize\bfseries, inner sep=2pt,
    fill=orange!12, draw=orange!60!black},
  shared/.style={draw, thick, rounded corners=3pt,
    minimum width=2.4cm, minimum height=0.6cm,
    font=\footnotesize, inner sep=2pt, fill=yellow!10,
    draw=yellow!60!black},
  flow/.style={-{Stealth[length=4pt, width=3pt]}, thick},
  biflow/.style={{Stealth[length=3pt]}-{Stealth[length=3pt]}, thick,
    blue!60!black},
  feedback/.style={{Stealth[length=3pt]}-{Stealth[length=3pt]}, thick,
    red!60!black, dashed},
  panellabel/.style={font=\small\bfseries, anchor=north west},
  levellabel/.style={font=\footnotesize\itshape, text=black!60},
]

% ── Panel positions ──────────────────────────────────────
\def\panelW{7.8}  % panel width
\def\panelH{3.8}  % panel height
\def\gapX{0.6}    % horizontal gap
\def\gapY{0.8}    % vertical gap

% Panel origins (bottom-left of each)
\coordinate (P0) at (0, \panelH + \gapY);
\coordinate (P1) at (\panelW + \gapX, \panelH + \gapY);
\coordinate (P2) at (0, 0);
\coordinate (P3) at (\panelW + \gapX, 0);

% ── Panel 0: Single Agent ───────────────────────────────
\begin{scope}[shift={(P0)}]
  \draw[black!20, rounded corners=6pt] (-0.3, -0.3)
    rectangle (\panelW + 0.3, \panelH + 0.3);
  \node[panellabel] at (0, \panelH + 0.15) {Level 0: Single Agent};
  \node[levellabel] at (\panelW, \panelH + 0.15) [anchor=north east]
    {No coordination};

  \node[agent, fill=blue!12, minimum width=2.2cm, minimum height=1.4cm]
    (a0) at (0.5*\panelW, 0.45*\panelH) {Agent};

  \node[font=\scriptsize, text=black!50, align=center]
    at (0.5*\panelW, 0.45*\panelH - 1.3)
    {All reasoning in one context.\\No dependencies to manage.};
\end{scope}

% ── Panel 1: Sequential Pipeline ────────────────────────
\begin{scope}[shift={(P1)}]
  \draw[black!20, rounded corners=6pt] (-0.3, -0.3)
    rectangle (\panelW + 0.3, \panelH + 0.3);
  \node[panellabel] at (0, \panelH + 0.15) {Level 1: Sequential Pipeline};
  \node[levellabel] at (\panelW, \panelH + 0.15) [anchor=north east]
    {Design-time coordination};

  \node[agent, fill=blue!12] (a1) at (1.2, 0.5*\panelH) {A};
  \node[agent, fill=blue!12] (b1) at (3.2, 0.5*\panelH) {B};
  \node[agent, fill=blue!12] (c1) at (5.2, 0.5*\panelH) {C};
  \node[agent, fill=green!12] (d1) at (7.0, 0.5*\panelH) {Out};

  \draw[flow] (a1) -- (b1);
  \draw[flow] (b1) -- (c1);
  \draw[flow] (c1) -- (d1);

  \node[font=\scriptsize, text=black!50, align=center]
    at (0.5*\panelW, 0.5*\panelH - 1.3)
    {Unidirectional flow. No feedback.\\Human architect resolves all coordination.};
\end{scope}

% ── Panel 2: Parallel Fork-Join ─────────────────────────
\begin{scope}[shift={(P2)}]
  \draw[black!20, rounded corners=6pt] (-0.3, -0.3)
    rectangle (\panelW + 0.3, \panelH + 0.3);
  \node[panellabel] at (0, \panelH + 0.15) {Level 2: Parallel Fork-Join};
  \node[levellabel] at (\panelW, \panelH + 0.15) [anchor=north east]
    {Orchestrator at merge};

  \node[agent, fill=purple!12] (orch2) at (1.0, 0.5*\panelH) {Orch};
  \node[agent, fill=blue!12] (w2a) at (3.5, 0.5*\panelH + 1.0) {B};
  \node[agent, fill=blue!12] (w2b) at (3.5, 0.5*\panelH) {C};
  \node[agent, fill=blue!12] (w2c) at (3.5, 0.5*\panelH - 1.0) {D};
  \node[agent, fill=purple!12] (merge2) at (6.0, 0.5*\panelH) {Merge};
  \node[agent, fill=green!12] (out2) at (7.4, 0.5*\panelH) {Out};

  \draw[flow] (orch2) -- (w2a);
  \draw[flow] (orch2) -- (w2b);
  \draw[flow] (orch2) -- (w2c);
  \draw[flow] (w2a) -- (merge2);
  \draw[flow] (w2b) -- (merge2);
  \draw[flow] (w2c) -- (merge2);
  \draw[flow] (merge2) -- (out2);

  % No inter-worker communication
  \draw[thick, red!40, dashed] (3.5, 0.5*\panelH + 0.5)
    -- node[font=\tiny, text=red!50, right, xshift=2pt] {no comm.}
    (3.5, 0.5*\panelH - 0.5);

  \node[font=\scriptsize, text=black!50, align=center]
    at (0.5*\panelW, -0.05)
    {Concurrent but isolated. Workers don't know each other exist.};
\end{scope}

% ── Panel 3: True MAS (Levels 3-4) ─────────────────────
\begin{scope}[shift={(P3)}]
  \draw[black!30, rounded corners=6pt, line width=1.2pt] (-0.3, -0.3)
    rectangle (\panelW + 0.3, \panelH + 0.3);
  \node[panellabel] at (0, \panelH + 0.15)
    {Levels 3--4: True Multi-Agent System};
  \node[levellabel, text=blue!60!black] at (\panelW, \panelH + 0.15)
    [anchor=north east] {Runtime coordination};

  % Shared state
  \node[shared] (ss) at (0.5*\panelW, 0.5*\panelH + 1.5)
    {Shared State};

  % Agents
  \node[agent, fill=blue!12] (ma) at (1.8, 0.5*\panelH) {A};
  \node[agent, fill=red!10] (mb) at (4.0, 0.5*\panelH + 0.6) {B};
  \node[agent, fill=blue!12] (mc) at (4.0, 0.5*\panelH - 0.6) {C};
  \node[agent, fill=red!10] (md) at (6.2, 0.5*\panelH) {D};

  % Human
  \node[human] (mh) at (0.5*\panelW, 0.5*\panelH - 1.5) {Human};

  % Bidirectional agent communication
  \draw[biflow] (ma) -- (mb);
  \draw[biflow] (ma) -- (mc);
  \draw[biflow] (mb) -- (md);
  \draw[biflow] (mc) -- (md);
  \draw[feedback] (mb) -- (mc);

  % Shared state connections
  \draw[flow, yellow!60!black, thin] (ma.north) -- (ss.south west);
  \draw[flow, yellow!60!black, thin] (mb.north) -- (ss.south);
  \draw[flow, yellow!60!black, thin] (md.north) -- (ss.south east);

  % Human connections
  \draw[biflow, orange!60!black] (mh) -- (ma);
  \draw[biflow, orange!60!black] (mh) -- (md);

  \node[font=\scriptsize, text=black!50, align=center]
    at (0.5*\panelW, -0.05)
    {Feedback loops, negotiation, shared state, self-interest, humans as peers.};
\end{scope}

\end{tikzpicture}
\caption{\textbf{The coordination complexity spectrum.}
  \emph{Level~0}: a single agent handles all reasoning in one context.
  \emph{Level~1}: agents execute sequentially with unidirectional flow; all
  coordination is resolved by the human architect at design time (adapted
  from~\citealt{cognition2025dontbuild}).
  \emph{Level~2}: subagents execute concurrently but cannot communicate with
  each other; the orchestrator handles coordination at merge (adapted
  from~\citealt{cognition2025dontbuild}).
  \emph{Levels~3--4}: true multi-agent systems with runtime dependencies,
  bidirectional communication (blue), feedback loops (red dashed), shared state
  (yellow), self-interested agents (red-tinted), and humans participating as
  peers (orange).  Classical MAS research studied Levels~3--4.  Most current
  LLM agent frameworks operate at Levels~1--2.}
\label{fig:coord-levels}
\end{figure*}


\begin{figure*}[t]
\centering
\begin{tikzpicture}[
  agent/.style={draw, thick, rounded corners=3pt, minimum width=1.1cm,
    minimum height=0.8cm, font=\footnotesize\bfseries, inner sep=2pt},
  human/.style={draw, thick, rounded corners=3pt, minimum width=1.1cm,
    minimum height=0.8cm, font=\footnotesize\bfseries, inner sep=2pt,
    fill=orange!12, draw=orange!60!black},
  shared/.style={draw, thick, rounded corners=3pt,
    minimum width=2.4cm, minimum height=0.6cm,
    font=\footnotesize, inner sep=2pt, fill=yellow!10,
    draw=yellow!60!black},
  flow/.style={-{Stealth[length=4pt, width=3pt]}, thick},
  biflow/.style={{Stealth[length=3pt]}-{Stealth[length=3pt]}, thick,
    blue!60!black},
  feedback/.style={{Stealth[length=3pt]}-{Stealth[length=3pt]}, thick,
    red!60!black, dashed},
  panellabel/.style={font=\small\bfseries, anchor=north west},
  levellabel/.style={font=\footnotesize\itshape, text=black!60},
]

% ── Panel positions ──────────────────────────────────────
\def\panelW{7.8}  % panel width
\def\panelH{3.8}  % panel height
\def\gapX{0.6}    % horizontal gap
\def\gapY{0.8}    % vertical gap

% Panel origins (bottom-left of each)
\coordinate (P0) at (0, \panelH + \gapY);
\coordinate (P1) at (\panelW + \gapX, \panelH + \gapY);
\coordinate (P2) at (0, 0);
\coordinate (P3) at (\panelW + \gapX, 0);

% ── Panel 0: Single Agent ───────────────────────────────
\begin{scope}[shift={(P0)}]
  \draw[black!20, rounded corners=6pt] (-0.3, -0.3)
    rectangle (\panelW + 0.3, \panelH + 0.3);
  \node[panellabel] at (0, \panelH + 0.15) {Level 0: Single Agent};
  \node[levellabel] at (\panelW, \panelH + 0.15) [anchor=north east]
    {No coordination};

  \node[agent, fill=blue!12, minimum width=2.2cm, minimum height=1.4cm]
    (a0) at (0.5*\panelW, 0.45*\panelH) {Agent};

  \node[font=\scriptsize, text=black!50, align=center]
    at (0.5*\panelW, 0.45*\panelH - 1.3)
    {All reasoning in one context.\\No dependencies to manage.};
\end{scope}

% ── Panel 1: Sequential Pipeline ────────────────────────
\begin{scope}[shift={(P1)}]
  \draw[black!20, rounded corners=6pt] (-0.3, -0.3)
    rectangle (\panelW + 0.3, \panelH + 0.3);
  \node[panellabel] at (0, \panelH + 0.15) {Level 1: Sequential Pipeline};
  \node[levellabel] at (\panelW, \panelH + 0.15) [anchor=north east]
    {Design-time coordination};

  \node[agent, fill=blue!12] (a1) at (1.2, 0.5*\panelH) {A};
  \node[agent, fill=blue!12] (b1) at (3.2, 0.5*\panelH) {B};
  \node[agent, fill=blue!12] (c1) at (5.2, 0.5*\panelH) {C};
  \node[agent, fill=green!12] (d1) at (7.0, 0.5*\panelH) {Out};

  \draw[flow] (a1) -- (b1);
  \draw[flow] (b1) -- (c1);
  \draw[flow] (c1) -- (d1);

  \node[font=\scriptsize, text=black!50, align=center]
    at (0.5*\panelW, 0.5*\panelH - 1.3)
    {Unidirectional flow. No feedback.\\Human architect resolves all coordination.};
\end{scope}

% ── Panel 2: Parallel Fork-Join ─────────────────────────
\begin{scope}[shift={(P2)}]
  \draw[black!20, rounded corners=6pt] (-0.3, -0.3)
    rectangle (\panelW + 0.3, \panelH + 0.3);
  \node[panellabel] at (0, \panelH + 0.15) {Level 2: Parallel Fork-Join};
  \node[levellabel] at (\panelW, \panelH + 0.15) [anchor=north east]
    {Orchestrator at merge};

  \node[agent, fill=purple!12] (orch2) at (1.0, 0.5*\panelH) {Orch};
  \node[agent, fill=blue!12] (w2a) at (3.5, 0.5*\panelH + 1.0) {B};
  \node[agent, fill=blue!12] (w2b) at (3.5, 0.5*\panelH) {C};
  \node[agent, fill=blue!12] (w2c) at (3.5, 0.5*\panelH - 1.0) {D};
  \node[agent, fill=purple!12] (merge2) at (6.0, 0.5*\panelH) {Merge};
  \node[agent, fill=green!12] (out2) at (7.4, 0.5*\panelH) {Out};

  \draw[flow] (orch2) -- (w2a);
  \draw[flow] (orch2) -- (w2b);
  \draw[flow] (orch2) -- (w2c);
  \draw[flow] (w2a) -- (merge2);
  \draw[flow] (w2b) -- (merge2);
  \draw[flow] (w2c) -- (merge2);
  \draw[flow] (merge2) -- (out2);

  % No inter-worker communication
  \draw[thick, red!40, dashed] (3.5, 0.5*\panelH + 0.5)
    -- node[font=\tiny, text=red!50, right, xshift=2pt] {no comm.}
    (3.5, 0.5*\panelH - 0.5);

  \node[font=\scriptsize, text=black!50, align=center]
    at (0.5*\panelW, -0.05)
    {Concurrent but isolated. Workers don't know each other exist.};
\end{scope}

% ── Panel 3: True MAS (Levels 3-4) ─────────────────────
\begin{scope}[shift={(P3)}]
  \draw[black!30, rounded corners=6pt, line width=1.2pt] (-0.3, -0.3)
    rectangle (\panelW + 0.3, \panelH + 0.3);
  \node[panellabel] at (0, \panelH + 0.15)
    {Levels 3--4: True Multi-Agent System};
  \node[levellabel, text=blue!60!black] at (\panelW, \panelH + 0.15)
    [anchor=north east] {Runtime coordination};

  % Shared state
  \node[shared] (ss) at (0.5*\panelW, 0.5*\panelH + 1.5)
    {Shared State};

  % Agents
  \node[agent, fill=blue!12] (ma) at (1.8, 0.5*\panelH) {A};
  \node[agent, fill=red!10] (mb) at (4.0, 0.5*\panelH + 0.6) {B};
  \node[agent, fill=blue!12] (mc) at (4.0, 0.5*\panelH - 0.6) {C};
  \node[agent, fill=red!10] (md) at (6.2, 0.5*\panelH) {D};

  % Human
  \node[human] (mh) at (0.5*\panelW, 0.5*\panelH - 1.5) {Human};

  % Bidirectional agent communication
  \draw[biflow] (ma) -- (mb);
  \draw[biflow] (ma) -- (mc);
  \draw[biflow] (mb) -- (md);
  \draw[biflow] (mc) -- (md);
  \draw[feedback] (mb) -- (mc);

  % Shared state connections
  \draw[flow, yellow!60!black, thin] (ma.north) -- (ss.south west);
  \draw[flow, yellow!60!black, thin] (mb.north) -- (ss.south);
  \draw[flow, yellow!60!black, thin] (md.north) -- (ss.south east);

  % Human connections
  \draw[biflow, orange!60!black] (mh) -- (ma);
  \draw[biflow, orange!60!black] (mh) -- (md);

  \node[font=\scriptsize, text=black!50, align=center]
    at (0.5*\panelW, -0.05)
    {Feedback loops, negotiation, shared state, self-interest, humans as peers.};
\end{scope}

\end{tikzpicture}
\caption{\textbf{The coordination complexity spectrum.}
  \emph{Level~0}: a single agent handles all reasoning in one context.
  \emph{Level~1}: agents execute sequentially with unidirectional flow; all
  coordination is resolved by the human architect at design time (adapted
  from~\citealt{cognition2025dontbuild}).
  \emph{Level~2}: subagents execute concurrently but cannot communicate with
  each other; the orchestrator handles coordination at merge (adapted
  from~\citealt{cognition2025dontbuild}).
  \emph{Levels~3--4}: true multi-agent systems with runtime dependencies,
  bidirectional communication (blue), feedback loops (red dashed), shared state
  (yellow), self-interested agents (red-tinted), and humans participating as
  peers (orange).  Classical MAS research studied Levels~3--4.  Most current
  LLM agent frameworks operate at Levels~1--2.}
\label{fig:coord-levels}
\end{figure*}
